\documentclass[pop,reprint]{revtex4-1}


\usepackage{fontspec}
\usepackage{indentfirst}
\usepackage{amsmath}
\usepackage{amsfonts}
\usepackage{amssymb}
\usepackage{bm}
\usepackage[bf,medium]{titlesec}
\usepackage{graphicx}% Include figure files

\setmainfont[BoldFont=cwTeXQHei-Bold]{cwTeX Q Ming}


\linespread{1.15}
\XeTeXlinebreaklocale "zh"
\XeTeXlinebreakskip = 0pt plus 1pt

%\titlelabel{\S\ \thetitle\quad}

%%%%%%%%%%%%%%%
%自定義的符號指令

\newcommand{\norm}[1]{\|#1\|}
\newcommand{\sdelta}{{\scriptstyle\Delta}}
\newcommand{\bra}[1]{\langle#1|}
\newcommand{\ket}[1]{|#1\rangle }
\newcommand{\braket}[2]{\langle#1|#2\rangle }
\newcommand{\proj}[2]{|#1\rangle\langle#2|}
\newcommand{\expect}[2]{\langle #1\rangle_{#2}}


%%%%%%%%%%%%%%%


\begin{document}

\title{Coherent State }
\author{Chih-Han Lin 林致翰}
\email{clin@ltl.iams.sinica.edu.tw}

\begin{abstract}

\end{abstract}
\maketitle
\section{相干態}
相干態（coherent state）即湮滅算子 $ \bm{a} $ \textbf{對應到的 eigenstate} $ \ket{a} $ 

\begin{eqnarray}
\bm{a}\ket{a}=a\ket{a}.
\end{eqnarray}

由於 $ \bm{a} $ 非 Heritian ，相干態的 eigenvalue 一般而言為複數。如果使用 $ \ket{n} $ 將 $ \ket{a} $ 展開，即 $ \ket{a}=\sum_n c_n\ket{n} $ ，$ c_n=\braket{n}{a} $ 則

\begin{eqnarray}
\bm{a}\ket{a}&=&\sum_{n=0}^\infty c_n\bm{a}\ket{n}=\sum_{n=1}^\infty c_n\sqrt{n}\ket{n-1}\\&=&a\sum_{n=0}^\infty c_n\ket{n}.
\end{eqnarray}

由此可知 $  c_n\sqrt{n}=ac_{n-1} $，故 $ c_1=ac_0 $、$ c_2=ac_1/\sqrt{2}=a^2c_0/\sqrt{2!} $、$ c_n=a^nc_0/\sqrt{n!} $，我們可以將 $ \ket{a} $ 改寫成

\begin{eqnarray}
\ket{a}=\sum_{n=0}^\infty c_n\ket{n}=c_0\sum_{n=0}^\infty\frac{a^n}{\sqrt{n!}}\ket{n}.
\end{eqnarray}

$ c_0 $ 可由歸一化條件 $ \braket{a}{a}=1 $ 來決定

\begin{eqnarray}
\protect\braket{a}{a}=1=|c_0|^2=\sum_{n=0}^\infty\frac{|a|^{2n}}{n!}\braket{n}{n}=|c_0|^2e^{|a|^2}.
\end{eqnarray}

因此 $ c_0 =e^{-|a|^2/2} $。 $ \ket{a} $ 投影到 $ \ket{n} $ 中可寫成

\begin{eqnarray}\label{coherent}
\ket{a}=e^{-|a|^2/2}\sum_{n=0}^\infty\frac{a^n}{\sqrt{n!}}\ket{n}.
\end{eqnarray} 

$ p(n)=|\braket{n}{a}|^2=e^{-|a|^2/2}/\sqrt{n!} $ ，相干態在 Fock state $ \ket{n} $ 的投影機率即 Poisson 分佈，$ p(n) $ 為出現 $ n $ 顆光子或激發態的機率。在 $ \ket{a} $ 中光子數的平均值即 $ \sum np(n)=|a|^2=\bra{a}\bm{a}^\dag\bm{a}\ket{a} $，其數目與 $ a $ 之 sqare norm 成正比。

$ \ket{a} $ 在 q-representation 中的波函數可寫成

\begin{eqnarray}
\notag\psi_a(x)&=&\braket{x}{a}=\braket{x}{n}\braket{n}{a}\\
\notag&=& \sqrt{\frac{1}{\sqrt{\pi}}}e^{-(|a|^2+x^2)/2}\sum_{n=0}^\infty\frac{(a/\sqrt{2})^n}{n!}H_n(x)\\
&=&\sqrt{\frac{1}{\sqrt{\pi}}}e^{-(x-\sqrt{2}x)^2/2}.
\end{eqnarray}

若將 $ a $ 寫為 $ a=x+iy $ ，我們可以將 $ a $ 表示成位置 $ x $ 與動量 $ p $ 在相干態下平均值的線性組合：

\begin{eqnarray}
\expect{x}{}&=&\bra{a}\bm{x}\ket{a}=\frac{1}{\sqrt{2}}\bra{a}(\bm{a}+\bm{a}^\dag)\ket{a}=\sqrt{2}x,\\
\expect{p}{}&=&\bra{a}\bm{p}\ket{a}=-\frac{(a-a^*)i\hbar}{\sqrt{2}}=\sqrt{2}\hbar y.
\end{eqnarray}

故 $ \displaystyle a=x+iy=\frac{1}{\sqrt{2}}\left(\expect{x}{}+\frac{1}{\hbar}\expect{p}{}\right) $. 相干態在 q-reprensentation 下的波函數可表達為平均值 $ \sqrt{2}a=\expect{x}{}+\expect{p}{}/\hbar $ 的 Guass 分佈。

\section{非正交性與超完備性}

\begin{eqnarray}
\notag\ket{a}\bra{a'}&=&\bra{n}\frac{e^{-|a|^2}a^{*n}}{\sqrt{n!}}\frac{e^{-|a'|^2}a'^{n'}}{\sqrt{n'!}}\ket{n'}\\
\notag&=&\frac{e^{-(|a|^2+|a'|^2)/2}(a^*a')^2}{n!}\\
\notag&=&\exp\left[-\frac{1}{2}(|a|^2+|a'|^2)+a^*a'\right]\\
&=&\exp\left(\frac{1}{2}|a-a'|^2\right).
\end{eqnarray}

由此可看出不同的相干態並不正交，只有在 $ a $ 與 $ a' $ 差距極大的時候才趨近正交。由相干態構成的投影算子 $ \ket{a}\bra{a} $ 是否具備完備性可進行複平面上的積分來計算：

\begin{eqnarray}
\notag\int dz\ket{z}\bra{z}&=&\sum_{n,n'}\frac{\ket{n}\bra{n'}}{\sqrt{n!}\sqrt{n'}!}
\int dz e^{-|a|^2}a^na^{*n'}\\
\notag&=&\sum_{n}\frac{\ket{n}\bra{n}}{n!}\int rdrd\theta e^{-r^2}r^{2n}\\
&=&2\pi\sum_{n}\frac{\ket{n}\bra{n}}{n!}\int_0^\infty dr r^{2n+1}e^{-r^2}=\pi
\end{eqnarray} 

我們可以把任意的 state $ \ket{\psi} $ 用相干態來展開

\begin{eqnarray}
\ket{\psi}=\frac{1}{\pi}\int da\ket{a}\braket{a}{\psi}.
\end{eqnarray}

由於相干態的非正交性，相干態的基底數超出空間維數，將任意的 $ \ket{\psi} $ 在相干態上展開時各個分量彼此線性相關，故謂之超完備性。

\section{相干態的平移算子}

欲找出 $ \bm{D}(a) $ 滿足 $  \bm{D}(a)\ket{a=0}=\ket{a} $ 可改寫 eq.~(\ref{coherent})

\begin{eqnarray}
\notag\ket{a}&=&e^{-|a|^2/2}\sum_{n=0}^\infty\frac{a^n}{\sqrt{n!}}\left(\frac{\bm{a}^{\dag n}}{\sqrt{n!}}\ket{n=0}\right)\\
\notag&=&e^{-|a|^2/2}\sum_{n=0}^\infty\frac{a^n\bm{a}^{\dag n}}{n!}\ket{n=0}\\
&=&e^{-|a|^2/2}e^{a\bm{a}^\dag}\ket{n=0}.
\end{eqnarray}

由於 $ e^{-a^*\bm{a}}\ket{n=0}=\ket{n=0} $ ，我們可以加入 $e^{-a^*\bm{a}}$ 因子到前一個式子中讓它看起來更對稱：

\begin{eqnarray}
\bm{D}(a)\ket{n=0}=e^{-|a|^2/2}e^{a\bm{a}^\dag}e^{-a^*\bm{a}}\ket{n=0}=\ket{a}.
\end{eqnarray}

利用 Campbell-Baker-Hausdorff 公式，已知

\begin{eqnarray}
\notag [a\bm{a}^\dag,[a\bm{a}^\dag,-a^*\bm{a}]]=0,\\
\notag [-a^*\bm{a},[-a^*\bm{a},a\bm{a}^\dag]]=0,\\
\protect[a\bm{a}^\dag, -a^*\bm{a}]=|a|^2.
\end{eqnarray}

可將位移算子寫成更精練的形式

\begin{eqnarray}
\bm{D}(a)=e^{-|a|^2/2}e^{a\bm{a}^\dag}e^{-a^*\bm{a}}=
e^{a\bm{a}^\dag-a^*\bm{a}}.
\end{eqnarray}

$ \bm{D}(a) $ 也是 Unitary 算子，$ \bm{D}(a)\bm{D}^\dag(a)=\bm{D}^\dag(a)\bm{D}(a)=\mathbf{I} $。$ \bm{D}(a) $ 有下述特性：

\begin{enumerate}
\item Unitary transform 形成複數位移 $ a $ 或 $ a^* $
\begin{eqnarray}
\bm{D}^\dag(a)F(\bm{a},\bm{a}^\dag)\bm{D}(a)=F(\bm{a}+a,\bm{a}^\dag+a^*).
\end{eqnarray}
\item 連續作用的位移算子與合成一次位移算子的效果僅有相位差
\begin{eqnarray}
\bm{D}(a)\bm{D}(a')=e^{(aa'^*-a^*a')/2}\bm{D}(a+a').
\end{eqnarray}
\item 不同位移算子彼此正交
\begin{eqnarray}
\mbox{Tr}[\bm{D}(a)\bm{D}^\dag(a')]=\pi\delta(\Re\{a-a'\})\delta(\Im\{a-a'\}).
\end{eqnarray}
\end{enumerate}
~\\

\subsection*{Reference}

\begin{footnotesize}
[1] L.Mandel, E.Wolf, \textsl{Optical Coherence and Quantum Optics}, Cambridge, 1995.

\noindent[2] 王正行, 量子力學原理, 2nd ed, 北京大學出版社, 2008.
\end{footnotesize}
\end{document}